\subsection{Odd terminal triples $(1,2,c)$ at a unit root}\label{proof:R2-C033}

\paragraph{R33.1 Exact object, direction and domains}

On the one original finite probability space let
\[
Jv=(Ev)1,\quad A=J+K,\quad T=J-K=A^*,\quad
g=K1,\quad p=1+g,\quad q=1-g,
\]
\[
\begin{gathered}
f_n=A^n1,\quad \bar f_n=T^n1,\quad h_n=Tf_{n-1},\quad\\
W_n=A\bar f_{n-1},\qquad I_q(V)=E[q^{-1/3}|V|^{8/3}].
\end{gathered}
\]
The terminal residual is exactly \((1,\allowbreak \allowbreak  2,\allowbreak \allowbreak  c)\), \(c\ge3\). The adjacent
internal root has exactly one residual unit arm. Its original inner
connector has length \(L\ge1\), and its original outer connector has
length two. Genuine equal pairs at either root are allowed with arbitrary
finite lengths and multiplicities.

The fixed actual message is
\[
H_D=E_r^-qA^L(E_t^-p\bar f_2W_c),\qquad U_D=T^2H_D.       \tag{R33.A1}
\]
The factors \(E_t^-,\allowbreak \allowbreak  E_r^-\) are the actual products of \(\bar f_\ell W_\ell\)
on original paired arms. The literal orientation is:


\begingroup\small
\setlength{\tabcolsep}{3pt}
\setlength{\LTleft}{0pt}\setlength{\LTright}{\fill}
\begin{longtable}{@{}p{0.2200\linewidth}p{0.7000\linewidth}@{}}
\caption{Odd terminal triples $(1,2,c)$ at a unit root: parameter values.}\label{tab:parameters:r2-c033:1}\\
\toprule
Entry & Direction/message \\
\midrule\endfirsthead
\toprule
Entry & Direction/message \\
\midrule\endhead
\bottomrule\endfoot
Terminal unit arm & Away from terminal, \(p\). \\
\addlinespace[2pt]
Terminal two-edge arm & Both edges toward terminal, \(\bar f_2\). \\
\addlinespace[2pt]
Terminal \(c\)-edge arm & First edge away from terminal; remaining edges toward its first vertex, \(W_c\). \\
\addlinespace[2pt]
Internal root unit arm & Toward the root, \(q\). \\
\addlinespace[2pt]
All \(L\) inner edges & From internal root toward terminal, \(A^L\). \\
\addlinespace[2pt]
Both outer edges & From internal root toward receiver, \(T^2\). \\
\addlinespace[2pt]
Each genuine equal pair & Entire reversal of its \(f_\ell h_\ell\) pair. \\
\end{longtable}
\endgroup


The parameters \(K,\allowbreak \allowbreak  \mu,\allowbreak \allowbreak  d\) and the spectral moments vary under this fixed rule.

The ordinary-fourth auxiliary is
\[
F_c=qf_2h_c,\qquad H_{c,L}=pT^LF_c.
\]
The exact intermediate assertion proved below is
\[
EH_{c,L}^4\le e^{-18d/5}
\quad\begin{cases}
0\le d\le1/3,&c=3,\\
0\le d\le1/4,&c=4,\\
0\le d\le1/5,&c\ge5,
\end{cases}
\quad L\ge1.                                      \tag{R33.A2}
\]
On the remaining high domains, compare \(W_c\) with one and apply
C026 in the same direction.

The final result is
\[
\boxed{I_q(U_D)^3\le e^{-197d/3420}\le1.}             \tag{R33.A3}
\]

\paragraph{R33.2 Original finite calculus and path means}

\paragraph{R33.2.1 Finite notation and the short spectral exception}
\label{part:R33.2.1}
Use $B=-K$, $u=-g$, $\mathsf X=-B^2$, $\tau=d-\delta_2$ and
\eqref{eq:odd-common:inputs}, with Lemma~\ref{lem:odd-common:spectral}.
Its original spectral measure $\nu$ at $u$ has mass $d$,
support $[0,\theta]$, and $\int x^j\,d\nu=\delta_{j+1}$.
Write $S_n^-=\one+G_n(B)u$ and
$\mu_n=\E f_n=\E\bar f_n$.
For $Q(B)=E_0(\mathsf X)+BD_0(\mathsf X)$ its two formulas are
\[
 E[Q(B)u]=-\int D_0\,d\nu,\qquad
 \|Q(B)u\|^2=\int(E_0^2+xD_0^2)\,d\nu.\tag{R33.A5}
\]
For $v=\delta_2/d$, Jensen gives
$v\ge d$, $\delta_3/d\ge v^2$ and $\delta_4/d\ge v^3$.
These are integrated moment restrictions, with no restriction $x\ge d$.
The positive convolution is \eqref{eq:Q188:conv}; the genuine
message and equal-pair bounds are \eqref{eq:odd-common:inputs}.
The spectral exception used below is
\[
 |G_n(i\sqrt x)|\le1\quad(n=1\text{ or }n\ge3),\qquad
 |G_2(i\sqrt x)|^2=1+x.\tag{R33.A7}
\]
The real coordinate for every $n\ge3$ lies in $[1-x,1]$.

\paragraph{R33.2.2 The all-length mean used in the high domain}
The stronger finite path estimate \eqref{eq:compact-strong-path-means}
gives, on the same original measure,
\[
 \boxed{\mu_n\le1-83d/32\quad(n\ge4).}\tag{R33.A9}
\]
Its proof includes the four short indices \(n=4,5,6,7\) and every
\(n\ge8\), with no monotonicity assumption on the integer \(n\).

\paragraph{R33.3 Signed-mean moment payments}
The common square payments
\eqref{eq:unit-two-step:square-log-scale} and
\eqref{eq:unit-two-step:variance-exponential} give, for every
real vector $V$,
\[
 Ee^{4BV}\le e^{4EBV+(17/4)\|V\|^2},\tag{R33.A13}
\]
\[
 Ee^{4BV}\le e^{4EBV+2\|V\|^2+6\|BV\|^2}.\tag{R33.A14}
\]
Both also hold with $K$ in place of $B$.

\paragraph{R33.4 Exact positive propagation and the ordinary-fourth outlet}

\paragraph{R33.4.1 The retained q square and all finite Duhamel coefficients}

Suppose a true nonnegative \(F\) obeys
\[
F\le q(m-U/2)^2,\qquad m>0,\quad k=1/m,\quad
V=EU^2,\quad R=-kU.
\]
Every application below obtains this square from the sum of two
nonnegative original arms; the square base is nonnegative.
Use the QS row of Table~\ref{tab:unit-two-step:parameters}:
$w=q$, $r=2$, $s_0=m^2$, $s_1=-m$, $s_2=1/4$ and $\lambda=1$.
The first row is
\[
 TF/m^2\le\bar f_2+BR+\delta,\qquad
 \delta=-k(EU-EgU)+k\sqrt{dV}+k^2V.\tag{R33.A15}
\]
Choose $D\ge\max(\delta,0)$ and put $D_i=EB^iR$.
The common finite continuation gives
\begin{align}
0\le T^LF/m^2\le Y_L
 &=\bar f_{L+1}+B^LR+D\bar f_{L-1}
       +\sum_{i=1}^{L-1}D_i\bar f_{L-1-i},\tag{R33.A16}\\
 Y_L&=S^-_{L+1}+B^LR+D\bar f_{L-1}
       -\sum_{j=3}^{L+1}\gamma_j\bar f_{L+1-j}.\tag{R33.A17}
\end{align}
Here $\gamma_j=\beta_j-D_{j-2}$. The index is $j-2$ because
the initial path is $\bar f_2$. Each local case below proves the
signs it uses before discarding a path term.

\paragraph{R33.4.2 The two fourth-moment payment formulas}

If $Y_L\le Cq+BPu+c$, $C>0$, apply
\eqref{eq:unit-two-step:unit-pairing} and the log-scale square
payment with scale $m^2C$:
\[
\begin{aligned}
 \log E(pT^LF)^4
 &\le8\log m+4\log C+\frac4C(c+EBPu)
                   +\frac{17}{4C^2}\|Pu\|^2.
\end{aligned}\tag{R33.A18}
\]
Increasing $C$ to a specified larger positive value keeps the
comparison because $q>0$. If instead $p+Y_L\le2+c+BV$,
the same payment with scale $m^2$ gives
\[
 \log E(pT^LF)^4
 \le8\log m+4(c+EBV)+(17/4)\|V\|^2.\tag{R33.A19}
\]

\paragraph{R33.4.3 The ordinary-fourth outlet with its exact curvature}
For an original positive kernel operator \(M\), row Cauchy and
H\"older with exponents \(3,3/2\) give
\[
 I_q(MY)^3\le\E[(M\one)^4/q]\,\|M\|^4(\E|Y|^4)^2.
 \tag{R33.A20}
\]
Apply the complete same-measure calculation in
paragraph~\ref{part:R29.11}. Its inverse row and ordinary norm use
\(\bar f_2=T^2\one\) and
\(\|T^2\|\le(1+1024\delta_2/855)^2\); the latter's Schur
proof is also given in Section~\ref{proof:R2-C026}, paragraph R26.8.1.
Thus, for every original real \(Y\),
\[
 I_q(T^2Y)^3\le B_2(\E|Y|^4)^2,\qquad
 B_2=[1-2d^2+(8-2d)\delta_2](1+1024\delta_2/855)^8,
\]
\[
 B_2\le\exp\{(24427/3420)d-5d^2\}.\tag{R33.A21}
\]
Consequently the ordinary-fourth input \(\E|Y|^4\le e^{-18d/5}\)
gives
\[
 I_q(T^2Y)^3\le e^{-(197/3420)d-5d^2}.\tag{R33.A22}
\]
The input is unweighted and real, and the outer operator and weight
remain the original \(T^2\) and \(q\).

\paragraph{R33.5 Terminal (1, 2, 3): every L and the whole original d domain}\label{part:R33.5}

\paragraph{R33.5.1 \(L=1:\) true positive square, actual moments and both continuous domains}

Let \(r=Kg,\allowbreak \allowbreak  \ s=K^2g,\allowbreak \allowbreak  \ t=1-d/2,\allowbreak \allowbreak  \ k=1/t\). The original paths are
\[
f_2=p+r,\qquad h_3=q-d-r-s,\qquad \bar f_3=q-d+r-s.
\]
Thus \(f_2+h_3=2-d-s\) and
\[
F_3\le q(t-s/2)^2.
\]
Apply Corollary~\ref{cor:unit-two-step:first-square} with
$m=t$, $a=0$, $W=s$, and $\Delta=\gamma$:
\[
 TF_3\le t^2q+tK(tg+s)+\gamma,\qquad
 \gamma=t(\sqrt{d\delta_3}-\delta_2)+\delta_3\ge0.
 \tag{R33.A23}
\]
Here $d\delta_3\ge\delta_2^2$. The resulting positive square is
\[
 H_{3,1}\le(c+N/2)^2,\qquad
 c=t+\gamma/(2t),\quad N=K(tg+s).\tag{R33.A24}
\]
The complete upper bound dominates $TF_3>0$; the corollary
therefore supplies the nonnegative square base.

The complete-square inequality
\(\sqrt{\theta v}-v\le\theta/4\), \(v=\delta_2/d\), yields
\[
c\le1-\left(\frac12-\frac{\theta}{8}
-\frac35\theta^2\right)d
=1-\frac{1793}{5120}d\le1-\frac7{20}d.
\]
For \(w=\delta_3/d\),
\[
EN=-(t-v)d,\quad
\|tg+s\|^2=d(t^2-2tv+w)\le d[t^2-(2t-\theta)v].
\]
Furthermore
\[
EN^2=\int x(t-x)^2\,d\nu\le dv(t-v)^2.            \tag{R33.A25}
\]
This is Jensen for the concave function \(x(t-x)^2\):
its second derivative is \(-4t+6x<0\) for
\(t\ge5/6,\allowbreak \allowbreak  \ x\le13/32\).

For \(v\le1/5\), (R33.A14) and the nonnegative base in (R33.A24) give
\[
\frac{\log EH_{3,1}^4}{d}
\le-\frac{14}{5}-4t+2t^2
+(4-4t+2\theta)v+6v(t-v)^2.
\]
The constant term is at most \(-427/90\);
\(4-4t+2\theta\le71/48\);
and \(v(1-v)^2\) is increasing on \([0,\allowbreak \allowbreak  1/5]\), with endpoint
\(16/125\). Therefore
\[
\begin{gathered}
\frac{\log EH_{3,1}^4}{d}\\
\le-\frac{427}{90}+\frac{71}{240}+\frac{96}{125}\\
=-\frac{18}{5}-\frac{1451}{18000}.
\end{gathered}\tag{R33.A26}
\]

For \(v\ge1/5\), retaining the exact \(\gamma\) and using (R33.A13) first
gives
\[
\begin{gathered}
-4-4t+\frac{17}{4}t^2+4\sqrt w\\
-\frac{17}{2}tv+\left(\frac4t+\frac{17}{4}\right)w.
\end{gathered}
\]
Only now use \(w\le\theta v\). The resulting \(\Phi(t,\allowbreak \allowbreak  v)\) is convex
in \(t\), with second derivative \(17/2+8\theta v/t^3>0\).
At its two endpoints:
\[
\Phi(1,v)=-15/4+4\sqrt{\theta v}-(659/128)v,
\]
\[
\Phi(5/6,v)=-631/144+4\sqrt{\theta v}-(6541/1920)v.
\]
Both decrease on \(v\ge1/5\), since
\(2\sqrt{\theta/v}<3\) while both subtracted coefficients exceed three.
The rational comparison \(\sqrt{13/160}<2/7\) follows from \(637<640\).
It gives
\[
\begin{gathered}
\Phi(1,1/5)<-\frac{18}{5}-\frac{33}{896},\\
\qquad\\
\Phi(5/6,1/5)<-\frac{31}{8}<-\frac{18}{5}.
\end{gathered}\tag{R33.A27}
\]
Thus every actual \(v\in[d,\allowbreak \allowbreak  \theta]\) is covered, with no spectral
support restriction beyond \([0,\allowbreak \allowbreak  \theta]\).

\paragraph{R33.5.2 All \(L\ge 2:\) the complete coefficient signs}

In (R33.A15) use
\[
m=1-d/2,\quad k=1/m,\quad U=s=\mathsf Xu,\quad
R=-k\mathsf Xu=kB^2u.
\]
Then \(EU-EgU=\delta_2\), \(EU^2=\delta_3\), and
\[
\begin{gathered}
\delta=-k\delta_2+k\sqrt{d\delta_3}+k^2\delta_3,\qquad\\
\frac{\delta}{d}\\
\le\frac{k\theta}{4(1-k\theta)}\\
\le\frac{39}{164}<\frac14.
\end{gathered}\tag{R33.A28}
\]
The first inequality completes the square in
\(k[\sqrt{\theta v}-(1-k\theta)v]\).
The exact finite coefficients are
\[
D_i=kEB^{i+3}1,\quad
D_1=k\delta_2,\ D_2=0,\ D_3=-k\delta_3,\ D_4=0.
\]
Every positive \(D_i\), \(i\ge3\), is at most
\(k\delta_4\le k\theta^3d<d/12\). Consequently
\[
\begin{gathered}
\gamma_3\ge41d/80,\quad \gamma_4=d,\quad\\
\gamma_j\ge49d/96=:\sigma\quad(j\ge5).
\end{gathered}\tag{R33.A29}
\]
The first bound exceeds \(\sigma\), since
\(41\cdot96-49\cdot80=16>0\).
For \(L=2\) I take \(D=d/4\); for all \(L\ge 3\) I take \(D=3d/8\).
Both dominate the actual \(\delta\).

\paragraph{R33.5.3 \(L=2\) in the canonical h3 direction}

The exact (R33.A16) reads
\[
Y_2=Cq+BQu+c,\quad
C=1+D,\quad c=k\delta_2-d,\quad Q=1+B+kB^3,\quad D=d/4.
\]
Thus
\[
c+EBQu=k\delta_2-2d\le-121d/80.
\]
At a frequency, \(Q=1+z(1-kx)\) and
\(0\le1-kx\le1-x\). Hence
\[
\|Qu\|^2\le d\left[1+\max_{0\le x\le1}x(1-x)^2\right]
=31d/27.
\]
Use \(C\le13/12\), \(\log m\le-d/2\), \(\log C\le d/4\)
in (R33.A18). The full payment is
\[
\begin{gathered}
\frac{\log EH_{3,2}^4}{d}\\
\le-4+1-\frac{363}{65}+\frac{527}{108}\\
=-\frac{26009}{7020}<-\frac{37}{10}<-\frac{18}{5}.
\end{gathered}\tag{R33.A30}
\]
This gives the canonical \(qf_2h_3\) fourth estimate.
The C106/C102 \(L=2\) eighth row concerns the direction
\(qf_2\bar f_3\).

\paragraph{R33.5.4 \(L=3\)}

The exact finite form is
\[
Y_3=Cq+BPu-d,\quad
C=1-d+D+k\delta_2\le\bar C=1-d/8,
\]
\[
P=1+D+B+(1-k\mathsf X)B^2,\qquad D=3d/8.
\]
After using \(d=2(1-1/k)\), the frequency difference
\((|P|^2-(1+D)^2)/x\) equals
\[
-5/2+3/(2k)+(7k/2-1/2)x-2kx^2+k^2x^3.
\]
It is convex in \(k\in[1,\allowbreak \allowbreak  6/5]\). Each endpoint is increasing in
\(x\in[0,\allowbreak \allowbreak  \theta]\), and its final value is respectively
\[
-1451/32768,\qquad -9507/204800.
\]
Thus \(\|Pu\|^2\le(1+D)^2d\).
The exact signed mean is
\[
-d+EBPu=-(2+D)d+\delta_2-k\delta_3\le-(7/4+D)d,
\]
using \(\delta_3/d\ge v^2\) and \(v-kv^2\le1/(4k)\le1/4\).
The ratio \((1+D)/\bar C\) increases to \(27/23<6/5\).
Since \(\bar C\le1\), division strengthens the negative mean, and
\[
\begin{gathered}
\frac{\log EH_{3,3}^4}{d}\\
\le-4-\frac12-7-4D+\frac{17}{4}\left(\frac65\right)^2\\
\le-\frac{269}{50}<-\frac{18}{5}.
\end{gathered}\tag{R33.A31}
\]

\paragraph{R33.5.5 \(L=4\)}

The actual finite expression is
\[
\begin{gathered}
Y_4=Cq+BPu+c,\quad C=1-2d+D+k\delta_2,\quad A_0=C+d,\\
P=A_0+(1+D)B+B^2+(1-k\mathsf X)B^3,\\
c=-\tau-Dd-k\delta_3,\qquad D=3d/8.
\end{gathered}
\]
At a frequency,
\[
|P|^2=A_0^2+x[(1+D-x+kx^2)^2+x-2A_0].
\]
The bracket is convex in \(x\); its second derivative is
\(2(-1+2kx)^2+4k(1+D-x+kx^2)>0\).
At zero it is at most
\(-1+2d+D^2\le-1/3+1/64<-1/4\).
At \(\theta\), use
\[
1+D-\theta+k\theta^2\le2347/2560<15/16,\quad A_0\ge19/24
\]
to obtain at most \(-229/768<-1/4\). Thus
\[
\|Pu\|^2\le A_0^2d-\delta_2/4.
\]
The legal enlarged coefficient is \(\bar C=1-9d/8\ge C\), and
\(A_0\le1-d/8\). The exact mean plus constant is
\[
c+EBPu=-2d+2\delta_2+d^2-2Dd-kd\delta_2-k\delta_3.
\]
Using the same \(v=\delta_2/d\), \(\delta_3/d\ge v^2\), \(v\ge d\) and
\(v\le\theta\), it is at most
\[
d[-19/16+d/4-2d^2]\le-53d/48.
\]
Also
\[
\|Pu\|^2/d\le1-d/2+d^2/64,\qquad
\frac{1-d/2+d^2/64}{(1-9d/8)^2}\le481/225.
\]
The last ratio increases, because the derivative numerator has the
sign of \(7/4-17d/32>0\), and its value at \(d=1/3\) is \(481/225\).
Therefore
\[
\begin{gathered}
\frac{\log EH_{3,4}^4}{d}\\
\le-4-\frac92-\frac{53}{12}+\frac{8177}{900}\\
=-\frac{862}{225}<-\frac{18}{5}.
\end{gathered}\tag{R33.A32}
\]

\paragraph{R33.5.6 Every finite \(L\ge 5,\allowbreak \) with no length cutoff}

The needed shifted finite-geometric assertion is
\[
\begin{gathered}
|G_n(i\sqrt x)-s_0(2+i\sqrt x)|\le1-2s_0\\
\quad(0\le s_0\le3/16,\ n\ge3).
\end{gathered}\tag{R33.A33}
\]
This finite two-term shift, including zero frequency and every
\(n\ge3\), is proved in paragraph~\ref{part:Q178.3.2}.

Retain the final four actual losses \(\bar f_0,\allowbreak \allowbreak  \bar f_1,\allowbreak \allowbreak  \bar f_2,\allowbreak \allowbreak  \bar f_3\) in (R33.A17),
using (R33.A29), and replace \(D\bar f_{L-1}\) by \(DS^-_{L-1}\).
The identity
\[
\bar f_0+\bar f_1+\bar f_2+\bar f_3=1+3q+B(2+B)u-d
\]
gives
\[
C=1+D-3\sigma=1-37d/32,\quad c=-\sigma(1-d),
\]
\[
P=G_L-k\mathsf XB^{L-1}+DG_{L-2}-\sigma(2+B).
\]
At every frequency \(0\le kx\le39/80\), and
\[
G_L-kxz^{L-1}=(1-kx)G_L+kxG_{L-1}.
\]
Apply (R33.A33) with \(s_0=\sigma=49d/96\le49/288<3/16\)
to both terms; \(G_{L-2}\) contracts because \(L-2\ge3\).
It follows that
\[
\|Pu\|^2\le a^2d,\qquad a=1-2\sigma+D=1-31d/48.
\]
The actual signed mean is
\[
EBPu=-\beta_{L+2}+D_L-D\beta_L+2\sigma d,
\]
so
\[
c+EBPu
\le d[-49/48+(335/256)d]\le-449d/768.
\]
Here \(-\beta_{L+2}+D_L=-\gamma_{L+2}\le-\sigma\);
every discarded loss has already been paid.
Finally \(0<C\le1\), \(a/C\le226/177<13/10\), and
\[
\begin{gathered}
\frac{\log EH_{3,L}^4}{d}\\
\le-4-\frac{37}{8}-\frac{449}{192}\\
+\frac{17}{4}\left(\frac{13}{10}\right)^2\\
=-\frac{18149}{4800}<-\frac{18}{5}.
\end{gathered}\tag{R33.A34}
\]
Equations (R33.A26)-(R33.A34) complete every \(c=3\) row on \(0<d\le1/3\).

\paragraph{R33.6 Terminal (1, 2, 4): every L on \(0<d\le 1/4\)}

\paragraph{R33.6.1 \(L=1\)}

Let \(r=Kg,\allowbreak \allowbreak  \ s=K^2g,\allowbreak \allowbreak  \ w=K^3g,\allowbreak \allowbreak  \ W=s+w\),
\(m=1-d,\allowbreak \allowbreak  \ k=1/m,\allowbreak \allowbreak  \ t=1-2d\).
The actual sum is
\[
f_2+h_4=2-2d+dg-W.
\]
Thus \(F_4\le q(m+U_0/2)^2\), \(U_0=dg-W\).
Put
\[
 \mathcal D=\theta(1+\theta)=585/1024<3/5,\qquad
 v=\delta_2/d,\quad y=\sqrt{\mathcal Dv}<1/2.
\]
The actual moments are $EW=\delta_2$,
$E(gU_0)=d^2+\delta_2$, and
$\|U_0\|\le\sqrt d(d+y)$.
Thus the q-square corollary permits
\[
 \Delta=\delta d,\qquad
 \delta=m(y-v)+(d+y)^2\ge0,
\]
since $\sqrt{d\|U_0\|^2}\le d(d+y)$.
For $N=K(tg+W)$ it gives directly
\[
 TF_4\le m^2q+mN-m\delta_2+\delta d,
 \qquad EN=-td+\delta_2,
\]
\[
 H_{4,1}\le(c_*+N/2)^2,\qquad
 c_*=m-\delta_2/2+\delta d/(2m),
\]
with the unchanged signed-vector bound
\[
 \|tg+W\|^2\le d[t^2-(2t-\mathcal D)v]\le dt^2.
\]
Its base is the arithmetic mean of $mp$ and the complete
positive upper factor.

Using (R33.A13), the two linear \(\delta_2\) terms cancel exactly:
\[
\begin{gathered}
\frac{\log EH_{4,1}^4}{d}\\
\le-8-4t+4(y-v)+4(d+y)^2/m\\
+\frac{17}{4}[t^2-(2t-\mathcal D)v].
\end{gathered}
\]
Because \(y-v\le\mathcal D/4<3/20\),
\(d+y<3/4\), \(m\ge3/4\), this is at most
\[
-22/5-4t+(17/4)t^2.
\]
This convex quadratic has maximum \(-83/20<-18/5\) on
\(t\in[1/2,\allowbreak \allowbreak  1]\). This proves the unit row without using the old
eighth bound as a fourth budget.

\paragraph{R33.6.2 The normalized input, quadratic error and every coefficient}

For \(L\ge 2\) use
\[
U=s+w-dg=(d+\mathsf X-\mathsf XB)u,\quad R=-kU,
\]
\[
EU^2=d^3+2d\delta_2+\delta_3+\delta_4,\qquad EU-EgU=d^2+2\delta_2.
\]
The symbols \(d^3\) and \(\delta_3\) are distinct.
Let \(J_0=EU^2/d\), \(v=\delta_2/d\). Then
\[
J_0\le d^2+[2d+\theta(1+\theta)]v.
\]
The scalar inequality
\(\sqrt{J_0}\le3/8+(2/3)J_0\), together with \(k\le4/3\), gives
\[
\delta/(kd)
\le3/8-d+2d^2+v[-2+4d+2\theta(1+\theta)]\le3/8.
\]
If the coefficient of \(v\) is nonpositive, discard it.
Otherwise use \(v\le\theta\); the remaining quadratic increases
in \(d\) and its value at \(d=1/4\) is
\(-1099/16384<0\). Thus \(D=d/2\) dominates \(\delta\).

For every \(j\ge1\), direct skew moments give
\[
D_{2j-1}=(-1)^{j+1}k(d\delta_j+\delta_{j+1}),\qquad
D_{2j}=(-1)^{j+1}k\delta_{j+2}.
\]
In particular
\[
D_1\le7d/8,\quad D_2<d/4,\quad
D_3=-k(d\delta_2+\delta_3),\quad D_4=-k\delta_4.
\]
The first later positive odd term is at most
\[
D_5\le(4/3)(1/4+\theta)\theta^2d
=1183d/8192<d/6,
\]
and \(D_6\le(4/3)\theta^4d<d/6\).
Each subsequent member of either positive class contracts by at
least \(\theta^2\) every four indices. Therefore
\[
\begin{gathered}
\gamma_3\ge d/8,\quad\gamma_4\ge3d/4,\quad\\
\gamma_j\ge41d/96=:\sigma\quad(j\ge5).
\end{gathered}\tag{R33.A35}
\]

\paragraph{R33.6.3 \(L=2\) and \(L=3\)}

For \(L=2,\allowbreak \)
\[
Y_2=Cq+BQu+c,\quad C=1+D,\quad
c=k(d^2+\delta_2)-d,
\]
\[
Q=(1-k\mathsf X^2)+[1-k(d+\mathsf X)]B.
\]
Both frequency coordinates are nonnegative:
\(k(d+x)\le7/8\) and \(kx^2\le169/768\).
Hence
\[
|Q|^2\le(1-x^2)^2+x(1-x)^2
=1+x-4x^2+x^3+x^4\le1+x-3x^2\le13/12.
\]
The actual mean is
\[
c+EBQu=-2d+k(d^2+\delta_2+\delta_3)\le-695d/768.
\]
Use \(C\le9/8,\allowbreak \allowbreak  \ \log m\le-d,\allowbreak \allowbreak  \ \log C\le d/2\) in (R33.A18):
\[
\begin{gathered}
\frac{\log EH_{4,2}^4}{d}\\
\le-8+2-\frac{695}{216}+\frac{221}{48}\\
=-\frac{1993}{432}<-\frac{18}{5}.
\end{gathered}\tag{R33.A36}
\]

For \(L=3,\allowbreak \)
\[
Y_3=Cq+BPu+c,\quad
C=1-d+D+D_1\le\bar C=1+3d/8,\quad c=-d+k\delta_3,
\]
\[
P=1+D+B+[1-k(d+\mathsf X)]B^2+k\mathsf XB^3.
\]
Its real coordinate is in \([0,\allowbreak \allowbreak  1+D]\), and its imaginary coefficient
is \(1-kx^2\in[0,\allowbreak \allowbreak  1]\).
Thus \(\|Pu\|^2\le d[(1+D)^2+\theta]\). The exact rational comparison
\[
\frac{(1+d/2)^2+\theta}{(1+3d/8)^2}\le45/32
\]
has positive cleared difference \(d(112-107d)/2048\).
The \(k\delta_3\) terms cancel in the signed mean:
\[
c+EBPu=-(2+D)d+(1-kd)\delta_2\le-51d/32.
\]
With \(\bar C\le35/32\),
\[
\begin{gathered}
\frac{\log EH_{4,3}^4}{d}\\
\le-8+\frac32-\frac{204}{35}+\frac{765}{128}\\
=-\frac{28457}{4480}<-\frac{18}{5}.
\end{gathered}\tag{R33.A37}
\]

\paragraph{R33.6.4 Every finite \(L\ge 4\)}\label{part:R33.6.4}

The finite filter is
\[
Q_L=G_L-k(d+\mathsf X)B^{L-1}+k\mathsf XB^L.
\]
At a frequency set \(a=k(d+x)\), \(b=kx\), so
\(0\le b\le a\le7/8\). Then
\[
\begin{gathered}
Q_L=(1-a)G_L+(a-b)G_{L-1}\\
+b[G_{L-1}+z^L].
\end{gathered}\tag{R33.A38}
\]
The coefficients are nonnegative and sum to one.
The extra finite polynomial \(H_L=G_{L-1}+z^L\) is
\(P_{L-1}\) of Section~\ref{proof:Q5-C033}.
For \(L\ge5\), the estimate following \eqref{eq:L33:finiteP} gives
\(|H_L|^2\le1-7x/15\le1\).
At \(L=4\), its squared modulus minus one is
\[
 x[-1+3x-2x^2+x^3]\le0;
\]
the bracket increases to \(-1451/32768\) at \(x=\theta\).
Thus \(|H_L|\le1\), hence \(|Q_L|\le1\), for every finite \(L\ge4\).

Retain the final \(\gamma_{L+1}\bar f_0\) loss in (R33.A17), and use
\(D\bar f_{L-1}\le DS^-_{L-1}\). The resulting true upper is
\[
Y_L\le(1+D)q+BPu-\gamma_{L+1},\qquad
P=Q_L+DG_{L-2}.
\]
By (R33.A35),
\[
EBPu=-\beta_{L+2}+D_L-D\beta_L\le-\sigma,
\quad c+EBPu\le-2\sigma.
\]
The exceptional \(G_2\) occurs only at \(L=4,\allowbreak \) and
\(\|G_2\|<6/5\). Uniformly,
\[
\frac{\|Pu\|}{(1+D)\sqrt d}
\le\frac{1+6D/5}{1+D}\le46/45.
\]
Substitution into (R33.A18) gives
\[
\begin{gathered}
\frac{\log EH_{4,L}^4}{d}\\
\le-8+2-\frac{82}{27}\\
+\frac{17}{4}\left(\frac{46}{45}\right)^2\\
=-\frac{9307}{2025}<-\frac{18}{5}.
\end{gathered}\tag{R33.A39}
\]
This completes the entire \(c=4\) low domain.

\paragraph{R33.7 Every \(c\ge 5,\allowbreak  L=1:\) the full continuous moment rectangle}

\paragraph{R33.7.1 Actual finite endpoints and positive auxiliary function}

Put
\[
\begin{gathered}
\alpha=\beta_c,\quad a=\beta_{c-1},\quad\\
m=1-(\alpha+a)/2,\quad k=1/m,\quad t=m-a,
\end{gathered}
\]
\[
\begin{gathered}
W=\sum_{j=3}^cK^j1=K^2G_{c-2}(K)g,\quad\\
U_0=ag-W,\quad\\
e_1=d-\beta_{c+1},\quad e_2=d-\beta_{c+2}.
\end{gathered}
\]
Here \(e_1,\allowbreak \allowbreak  e_2\) denote the short means called \(D_1,\allowbreak \allowbreak  D_2\) in
R2-C033 \mbox{\hyperref[part:R33.5]{part R33.5}}; they are not the later Duhamel sequence.
The actual positive convolution, retaining \(aq\), gives
\[
\begin{gathered}
h_c\le(1-a)q-r-W-\alpha,\qquad\\
f_2+h_c\le2m+U_0,\qquad\\
F_c\le q(m+U_0/2)^2.
\end{gathered}\tag{R33.A40}
\]

The first-connector square now follows from
Corollary~\ref{cor:unit-two-step:first-square}; its actual moments
are evaluated next.

\paragraph{R33.7.2 Same-moment errors, finite spectral loss and endpoint coupling}\label{part:R33.7.2}

Let \(v=\delta_2/d\in[d,\allowbreak \allowbreak  \theta]\), \(y=\sqrt{\theta v}\).
Since \(c-2\ge3\), (R33.A7) gives
\[
\|U_0\|\le\sqrt d(a+y),\qquad
E(gU_0)=ad+e_2\ge(a+\kappa v)d.
\]
The same expansion as in the c4 unit row therefore permits
\[
\delta=m(y-\kappa v)+(a+y)^2\ge0,
\]
\[
N=K(tg+W),\quad TF_c\le m^2q+mN-me_1+\delta d,\quad EN=-td+e_2,
\]
\[
\begin{gathered}
H_{c,1}\le(c_*+N/2)^2,\quad\\
c_*=m-e_1/2+\delta d/(2m).
\end{gathered}\tag{R33.A42}
\]
Again the square is a nonnegative arithmetic mean of the actual
positive factors.

For every finite \(n\ge3\),
\[
\begin{gathered}
|G_n(i\sqrt x)|^2\\
\le\frac{(1+x^3)^2}{1+x}\le1-\eta x,\qquad\eta=7/16.
\end{gathered}\tag{R33.A43}
\]
The first bound follows separately from the three finite parity
formulas in \mbox{\hyperref[part:R33.2.1]{part R33.2.1}}.
For the second, clear \(1+x\): the remaining inequality is
\(9-7x-32x^2-16x^5\ge0\).
It decreases on \([0,\allowbreak \allowbreak  \theta]\), and
\[
9-7\theta-32\theta^2=7/8,\qquad16\theta^5<1/2.
\]
Thus it holds on the complete continuous interval.
Using also \(\Re G_n\ge1-x\),
\[
|t-xG_n|^2\le t^2-2tx+(2t+1)x^2-\eta x^3.
\]
Integration with the same \(\nu\), then \(\delta_3\le\theta \delta_2\) and
\(\delta_4/d\ge v^3\), gives
\[
\begin{gathered}
\|tg+W\|^2\le dA(t,v),\qquad\\
A=t^2-(2\kappa t-\theta)v-\eta v^3.
\end{gathered}\tag{R33.A44}
\]
Here \(m\ge4/5,\allowbreak \allowbreak  \ t\ge3/5\), so \(2\kappa t-\theta>0\).
The bound is positive on the entire enlarged rectangle:
it increases in \(t\), decreases in \(v\), and at
\(t=3/5,\allowbreak \allowbreak  v=\theta\) exceeds
\[
9/25-1/8-3/100=41/200.
\]

The short endpoint difference must be coupled to \(\alpha\):
\[
e_2-e_1\le h(\alpha-\tau),\qquad h=\theta/\kappa=13/19. \tag{R33.A45}
\]
For \(c=5,\allowbreak \) 6, \(\alpha=\tau\) and the difference is nonpositive.
For \(c\ge 7,\allowbreak \) the alternating endpoints give \(e_2-e_1\le \delta_4\),
while \(\alpha-\tau\ge\kappa \delta_3\) and \(\delta_4\le\theta \delta_3\).
This verifies (R33.A45) for every finite c and preserves the same-kernel
relation rather than separately maximizing its two terms.

\paragraph{R33.7.3 Monotonicity in both structural parameters and the whole v interval}

By (R33.A13) and (R33.A42)-(R33.A45), the logarithmic fourth bound divided by d
is at most
\[
\begin{gathered}
\mathcal E=\\
-4(\alpha+a)/d-4t+4h(\alpha-\tau)/d\\
+4(y-\kappa v)+4(a+y)^2/m+(17/4)A(t,v).
\end{gathered}\tag{R33.A46}
\]
Enlarge to the full rectangle \(\alpha,\allowbreak \allowbreak  a\in[\tau,\allowbreak \allowbreak  d]\), with the
original \(d,\allowbreak \allowbreak  v\) fixed.
There \(m\ge4/5,\allowbreak \allowbreak  \ t\ge3/5,\allowbreak \allowbreak  \ (a+y)/m<4/5\), and \(t-\kappa v>0\).
Direct differentiation gives
\[
\mathcal E_\alpha
=-4(1-h)/d+2+2[(a+y)/m]^2-(17/4)(t-\kappa v)<0,
\]
\[
\mathcal E_a
=-4/d+6+8(a+y)/m+2[(a+y)/m]^2-(51/4)(t-\kappa v)<0.
\]
The first is bounded above, after discarding its negative final term, by
\(-120/19+2+32/25<0\); the second by
\(-20+6+32/5+32/25<0\).
Hence the single upper bound (R33.A46) is maximized at
\(\alpha=a=\tau=(1-v)d\).

Now
\[
m=1-(1-v)d,\quad t=1-2(1-v)d,\quad z=(1-v)d+y.
\]
With \(y'=\frac12\sqrt{\theta/v}\ge1/2\),
\[
(z^2/m)'=(z/m)[2y'-2d-dz/m]>0.
\]
Also
\[
A'=4d(t-\kappa v)-(2\kappa t-\theta)-3\eta v^2>-1.
\]
Indeed \(2\kappa-\theta=25/32\) and
\(3\eta\theta^2<7/32\). The total derivative of (R33.A46) is therefore
bounded below by
\[
8-\frac85-\frac38-\frac{17}{4}=\frac{71}{40}>0.
\]
This pays the entire \(v\) interval, so it is enough to set
\(v=\theta,\allowbreak \allowbreak  \ \alpha=a=\kappa d,\allowbreak \allowbreak  \ y=\theta\).

\paragraph{R33.7.4 The final continuous d interval}

At those endpoints, (R33.A44) becomes
\[
P(d)=342509/524288-(14763/8192)d+(361/256)d^2.
\]
The remaining bound is exactly
\[
\begin{gathered}
\mathcal E_*(d)=\\
-8\kappa-4(1-2\kappa d)+4\theta^2\\
+\frac{4(\kappa d+\theta)^2}{1-\kappa d}\\
+(17/4)P(d).
\end{gathered}\tag{R33.A47}
\]
It is convex on \([0,\allowbreak \allowbreak  1/5]\): \(P''>0\), and the square of an affine
function divided by a positive affine function has nonnegative
second derivative here.
At zero, \(P(0)<2/3\), \(8\theta^2<4/3\), so
\[
\mathcal E_*(0)<-35/4+25/6=-55/12.
\]
At \(d=1/5\),
\[
t=61/80,\quad m=141/160,\quad z=21/40,\quad
P(1/5)=4577893/13107200<7/20,
\]
\[
4z^2/m=294/235<13/10,\qquad4\theta^2<2/3.
\]
Thus
\[
\begin{gathered}
\mathcal E_*(1/5)\\
<-39/5+2/3+13/10+119/80\\
=-1043/240<-18/5.
\end{gathered}\tag{R33.A48}
\]
Convexity proves the entire d interval. Together with the preceding
monotonicities and finite-tail arguments, this completes every
\(c\ge 5\) unit-inner row in the exact low domain.

\paragraph{R33.8 Every \(c\ge 5,\allowbreak  L\ge 2:\) all retained masks and signed propagation}

Throughout this section \(0<d\le1/5\), \(m=1-9d/8\), and
\(k=1/m\le40/31\).

\paragraph{R33.8.1 Genuine retained coefficients for every finite c}
Use Lemma~\ref{lem:unit-retained-pair} with $n=c$ and $\sigma=9/4$.
Its genuine coefficients $a_1,a_2,a_3$, sums $C_0,D_0,E_0$, and
actual vector $U$ are used here. It includes $c=5$ with $a_3=0$,
and every $c\ge6$ with $0\le\lambda\le1$. The second actual arm sum gives
\[
 f_2+h_c\le2m-U,\qquad F_c\le q(m-U/2)^2.\tag{R33.A49}
\]
Its exact constant loss and coefficient ranges are
\[
 \beta_c+a_1+a_2+(1-d)a_3=9d/4,
\]
\[
 C_0=9d/4-\beta_c+dE_0,\quad
 D_0=9d/4-\beta_c-\beta_{c-1}+dE_0,
\]
\[
 \begin{gathered}
 5d/4+dE_0\le C_0\le(53/32)d+dE_0,\\
 d/4\le D_0\le(17/16+d)d,\qquad0\le E_0\le d.
 \end{gathered}\tag{R33.A50}
\]
The ranges follow from $\tau\le\beta_j\le d$ and $E_0\le d$.

\paragraph{R33.8.2 Every finite frequency mask, including $c=5$}
Write $U=[a(\mathsf X)-b(\mathsf X)B]u$. The same lemma gives
\[
 a=5d/4+\delta_2+x-x^2,\qquad b=d/4+\delta_2+x
 \quad(c=5),\tag{R33.A51}
\]
and, for $c\ge6$, with $G_{c-2}(i\sqrt x)=a_f+i\sqrt x\,b_f$,
\[
 a=C_0+x(a_f-E_0),\qquad b=D_0+xb_f.\tag{R33.A52}
\]
The finite alternating sums obey
$1-x\le a_f,b_f\le1-x+x^2$ and $a_f\ge(1-x)b_f$. In both cases,
\[
 a,b\ge0,\qquad a-(1-x)b\ge a_1+a_2x\ge\tau.\tag{R33.A53}
\]
The following two paragraphs prove the full-frequency bounds and the
quadratic error for these same actual coefficients, with the original
global $\delta_2$ fixed.

\paragraph{R33.8.3 Strong retained norms}
Lemma~\ref{lem:reference-retained:strong-norm} applies to this same actual
\(U=[a(\mathsf X)-b(\mathsf X)B]u\), with the original global
\(d,\delta_2\) fixed during every frequency maximization. It gives
\[
 \sqrt{a^2+xb^2}\le u_0:=3/8+9d/5\le147/200,
 \qquad b\le z_0:=\theta+13d/16\le91/160.\tag{R33.A54}
\]

\paragraph{R33.8.4 The retained first-row allowance}\label{part:R33.8.4}
The \(\sigma=9/4\) row of Lemma~\ref{lem:reference-retained-packet}
proves, for these actual coordinates and the whole frequency interval,
\[
 \frac{15}{7}(a^2+xb^2)\le a+b.\tag{R33.A57}
\]
With \(V=\E U^2\), the actual first-row error is
\[
 \delta=-k\int(a+b)\,d\nu+k\sqrt{dV}+k^2V
       \le3kd/10\le12d/31<2d/5.
\]
Thus the same positive continuation uses
\[
 D=2d/5.\tag{R33.A59}
\]

\paragraph{R33.8.5 Signed finite path losses}
Use the exact four sign classes in
\eqref{eq:reference-retained:signed-moments} for this same \(R=-kU\).
The coordinate bounds of that lemma imply the original strong estimates
\[
 D_1\le(147/155)d<d,\quad D_2\le(1183/3968)d<d/3,
 \quad D_1+D_2\le(24731/19840)d<5d/4.\tag{R33.A61}
\]
Every subsequent positive class is below \(d/6\), so
\[
 \gamma_3>0,\quad\gamma_4\ge2d/3,\quad
 \gamma_j\ge41d/96\quad(j\ge5).\tag{R33.A62}
\]

\paragraph{R33.8.6 Every finite \(L\ge 4\)}

The filter is
\[
Q_L=G_L-kaB^{L-1}+kbB^L.
\]
At each frequency \(0\le ka\le147/155<1\) and
\(0\le kb\le91/124<1\). It is exactly the convex combination of
\[
G_{L-1},\quad G_L,\quad G_{L-1}+z^L,\quad G_{L+1}
\]
with respective weights
\[
(ka)(1-kb),\quad(1-ka)(1-kb),\quad(ka)(kb),\quad(1-ka)(kb).
\]
The four weights sum to one. All four finite polynomials contract:
the G polynomials by (R33.A7), the missing-term polynomial by the
complete proof in \mbox{\hyperref[part:R33.6.4]{part R33.6.4}}.
Therefore \(\|Q_Lu\|^2\le d\) for every finite \(L\ge 4.\)

Retain the final loss in (R33.A17) and use \(D\bar f_{L-1}\le2D\):
\[
p+Y_L\le2+c+BQ_Lu,\qquad c=2D-\gamma_{L+1}.
\]
The exact mean is
\[
EBQ_Lu=-\beta_{L+2}+D_L=-\gamma_{L+2}\le-\sigma.
\]
Hence
\[
c+EBQ_Lu\le4d/5-82d/96=-13d/240.
\]
The right side is a true nonnegative upper sum. Since
\(8\log m\le-9d\), (R33.A19) gives
\[
\begin{gathered}
\frac{\log EH_{c,L}^4}{d}\\
\le-9-\frac{13}{60}+\frac{17}{4}\\
=-\frac{149}{30}<-\frac{18}{5}.
\end{gathered}\tag{R33.A63}
\]

\paragraph{R33.8.7 The two short connectors}
The retained bounds give \(0\le ka\le147/155<1\) and
\(0\le kb\le91/124<1\). At \(z=i\sqrt x\), the two filters are
\[
 Q_2=1-kbx+z(1-ka),\qquad
 Q_3=1-x+kax+z(1-kbx).
\]
Their real and normalized imaginary coordinates lie in \([0,1]\).
Consequently
\[
 |Q_j|^2\le1+x\le45/32,\qquad
 \|Q_ju\|^2\le(45/32)d\quad(j=2,3).                 \tag{R33.A64}
\]

For \(L=2\), the original finite expansion is
\[
\begin{gathered}
 p+Y_2=2+c+BV,\qquad
 c=D+D_1-d,\qquad V=Q_2u+D\,1_{\rm const},\\
 EBV=-d+D_2,\qquad EQ_2u=-d+D_1\le0 .
\end{gathered}                                                     \tag{R33.A65}
\]
Thus \(c+EBV\le-7d/20\), and the constant-vector cross term is
nonpositive:
\[
 \frac{\|V\|^2}{d}\le\frac{45}{32}+\frac4{25}d
 \le\frac{5753}{4000}<\frac32.
\]
Here and below \(D=2d/5\). Formula (R33.A19), with
\(8\log m\le-9d\), gives
\[
 \frac{\log EH_{c,2}^4}{d}
 \le-9-\frac75+\frac{51}{8}
 =-\frac{161}{40}<-\frac{18}{5}.                    \tag{R33.A66}
\]

\paragraph{R33.8.8 The connector \(L=3\)}
The exact expansion, including its constant vector, is
\[
\begin{gathered}
 p+Y_3=2+c+BV,\quad
 \lambda_0=d-D-D_1,\quad P=Q_3+D,\\
 c=D+D_1+D_2-2d,\qquad V=Pu-\lambda_0 1_{\rm const},\\
 EBV=-\tau+D_3-Dd,\qquad EPu=-d+D_2 .
\end{gathered}                                                     \tag{R33.A70}
\]
The signed-moment bounds give
\[
 c+EBV\le[2/5+5/4-2-\kappa]d
 =-151d/160.                                         \tag{R33.A71}
\]
Since \(0\le D_1<d\) and \(0\le D_2<d/3\),
\( |\lambda_0|\le3d/5\) and \(|EPu|\le d\).
The constant-vector correction therefore satisfies
\[
 -2\lambda_0 EPu+\lambda_0^2\le(39/25)d^2.
\]
Using \(\sqrt{45/32}<6/5\) for the additional \(Du\) term gives
\[
\begin{aligned}
 \|V\|^2/d
 &\le45/32+(24/25)d+(4/25)d^2+(39/25)d\\
 &\le38333/20000<2.
\end{aligned}
\]
Finally (R33.A19) gives
\[
 \frac{\log EH_{c,3}^4}{d}
 \le-9-\frac{151}{40}+\frac{17}{2}
 =-\frac{171}{40}<-\frac{18}{5}.                    \tag{R33.A73}
\]
Together with part R33.7 and the preceding long rows, these two
estimates prove every \(c\ge5\) low-domain case.

\paragraph{R33.9 The same actual direction on the low and high domains}

\paragraph{R33.9.1 Low-domain inner reversal followed by the original outer \(T^2\)}\label{part:R33.9.1}

parts R33.5-8 prove (R33.A2) for the actual canonical input.
Apply that statement to the legal kernel \(-K\). Then
\[
p\leftrightarrow q,\quad A\leftrightarrow T,\quad
f_2\mapsto \bar f_2,\quad h_c\mapsto W_c,
\]
and all original moments \(\delta_i\) stay unchanged. Thus
\[
E[qA^L(p\bar f_2W_c)]^4\le e^{-18d/5}                 \tag{R33.A74}
\]
in precisely the stated low domains.
This step reverses the inner component only.
The actual output is then obtained by applying the
original \(T^2\), with the original receiver weight q.
For the genuine decorations,
\[
0<H_D=E_r^-qA^L(E_t^-p\bar f_2W_c)\le qA^L(p\bar f_2W_c),
\]
because \(0<E_t^-,\allowbreak \allowbreak  E_r^-\le1\) and \(A^L\) is positive.
Hence \(EH_D^4\le e^{-18d/5}\).
Apply the original ordinary-fourth bridge (R33.A21)-(R33.A22) directly to
this actual \(H_D\):
\[
I_q(U_D)^3\le e^{-(197/3420)d-5d^2}.              \tag{R33.A75}
\]
This proves all \(c=3\) hosts, the \(c=4\) hosts with \(d\le 1/4,\allowbreak \) and all \(c\ge 5\)
hosts with \(d\le 1/5.\)

\paragraph{R33.9.2 Actual high-domain path comparison}

For \(c=4\) and \(d\ge 1/4,\allowbreak \)
\[
W_4=A\bar f_3\le2E\bar f_3=2(1-2d)\le1.
\]
For \(c\ge 5\) and \(d\ge 1/5,\allowbreak \) the all-finite mean bound (R33.A9) gives
\[
\begin{gathered}
W_c=A\bar f_{c-1}\le2\mu_{c-1}\\
\le2(1-83d/32)\le77/80<1.
\end{gathered}\tag{R33.A76}
\]
Consequently the very same direction satisfies
\[
0<U_D\le U_{12,L}:=T^2[qA^L(p\bar f_2)].              \tag{R33.A77}
\]
This is a comparison inside a nonnegative integral. The actual
long arm and all its original edges have not been removed from
the oriented object or its final edge count.

The right side of (R33.A77) is the actual terminal $(1,2)$ direction of Section~\ref{proof:R2-C026}, with the same original $q$ and outer $T^2$.

The high-domain comparison table is:

\begin{description}[style=nextline,leftmargin=1.6em]
\item[$L=1$.] $\E[U_{12,L}^4/q]\le1-d/5$, by R26.3.
\item[$L=2$.] $\E[U_{12,L}^4/q]\le1-d/24$, by R26.4.2.
\item[$L=3$.] $\E[U_{12,L}^4/q]\le1-7d/10$, by R26.4.3.
\item[$L=4$.] $\E[U_{12,L}^4/q]\le1-d/4$, by R26.4.4.
\item[$L\ge5$.] $\E[U_{12,L}^4/q]\le1-d/2$, by R26.6.
\end{description}

The proofs at $L=3,4$ propagate constants as $T^2\one=\bar f_2$. The proof for every $L\ge5$ retains its finite final monomial and the covariance sign on the spectral measure of $\one$. Only these five actual terminal $(1,2)$ entries are used here.

All entries are at most \(1-d/24\). Positivity, \(q>0\) and (R33.A77)
therefore give
\[
E(U_D^4/q)\le1-d/24.
\]
Hölder with exponents \(3/2,\allowbreak \allowbreak  3\), using \(Eq=1\), is exactly
\[
I_q(U_D)
=E[(U_D^4/q)^{2/3}q^{1/3}]
\le[E(U_D^4/q)]^{2/3}.
\]
Thus the high-domain payment is
\[
\begin{gathered}
I_q(U_D)^3\\
\le(1-d/24)^2\le e^{-d/12}\\
\le e^{-197d/3420}.
\end{gathered}\tag{R33.A78}
\]
The final coefficient margin is
\(1/12-197/3420=22/855>0\).
Both split endpoints \(d=1/4\) and \(d=1/5\) belong to both domains.
Combining (R33.A75) and (R33.A78) proves (R33.A3).

\paragraph{R33.10 Original decorations, reversal, KL and exact gluing responsibilities}

\paragraph{R33.10.1 Original edges and all repeated arms}

Each genuine pair is formed from two actually present equal-length
arms. The pointwise inequality
\[
\bar f_\ell W_\ell\le1
\]
is the complete reversal of \(f_\ell h_\ell\le1\).
It permits comparison at the same root label, followed by positive
propagation.
The terminal pair factor remains inside \(A^L\), and the internal-root
pair factor remains outside \(A^L\) and before \(T^2\).

If repeated arms need a tie-breaking choice, any fixed labeling or
structure-based choice may be made before the host is supplied.
Every original arm in each pair is still oriented, and every original
subdivision vertex is still integrated over the original \(\mu\).


Without paired decorations the component has
\[
e_B=(1+2+c)+L+1+2=L+c+6
\]
original edges. If \(S_{\rm pair}\) is the sum of one arm length from
each actual pair at both roots, then
\[
e_B=L+c+6+2S_{\rm pair},\qquad |V_B|-1=e_B.         \tag{R33.A79}
\]
Every two-edge path, every internal connector edge and every long-arm
subdivision edge appears exactly once in (R33.A1).

\paragraph{R33.10.2 Complete component reversal is a separate operation}

Reversing every original edge means substituting \(-K\) throughout
the complete component, including both outer edges and every
decoration. The fully reversed output is
\[
U_D^{\rm rev}
=A^2[E_r^+pT^L(E_t^+qf_2h_c)],
\]
and its budget is
\[
I_p(U_D^{\rm rev})^3\le e^{-197d/3420},\qquad
I_p(V)=E[p^{-1/3}|V|^{8/3}].
\]
This is different from the inner reversal in \mbox{\hyperref[part:R33.9.1]{part R33.9.1}}.
The outer operator and the receiver weight must both swap in a
complete reversal. Keeping \(T^2\) while changing only the weight,
or keeping q while changing only the outer path, is not approved.
